\subsection{三角関数による表示}

$x^{p-1}(1-x)^{q-1}$の$x$と$1-x$は,和が$1$である.和が$1$になる2つの量といえば$\sin^2t$と$\cos^2t$であろう.そこで$x=\sin^2t$と置くと,\,$1-x=\cos^2t$であり,\,$0<t<\frac{\pi}{2}$で
\[
 \dd x=2\sin t\cos t\dd t
\]
である.この範囲では$\sin t,\cos t$が正なので,実数の指数もそのまま扱える.式\eqref{式:ベータ関数の定義}は
\begin{equation}\label{式:ベータ三角表示}
 B(p,q)=2\int_0^{\frac{\pi}{2}}
 (\sin t)^{2p-1}(\cos t)^{2q-1}\dd t
\end{equation}
となる.端点で関数が定まらない場合は,端点を避けて置換した後で極限を取る.左辺の両端の積分が収束しているので,対応する右辺の極限も有限である.

$p=q=\frac12$なら両方の指数が$0$になり,
\[
 B\left(\frac12,\frac12\right)
 =2\int_0^{\frac{\pi}{2}}1\dd t=\pi
\]
である.つまり
\[
 \int_0^1\frac1{\sqrt{x(1-x)}}\dd x=\pi
\]
となる.両端で高さが限りなく大きくなるにもかかわらず,積分は有限になる具体例である.

一方,\,$p=\frac{n+1}{2}$,\,$q=\frac12$とすれば,\,$\cos$の指数が$0$になって
\[
 B\left(\frac{n+1}{2},\frac12\right)=2\int_0^{\frac\pi2}\sin^nt\dd t=2I_n
\]
となる.\,$p$を$1$増やすことは,\,$\sin$の指数$n$を$2$増やすことに対応する.実際,式\eqref{式:ベータ漸化式}で$p=\frac{n+1}2$,\,$q=\frac12$とすると,\,$\frac{p}{p+q}=\frac{n+1}{n+2}$なので
\[
 B\left(\frac{n+3}{2},\frac12\right)
 =\frac{n+1}{n+2}
  B\left(\frac{n+1}{2},\frac12\right)
\]
となる.両辺を$2$で割れば$I_{n+2}=\frac{n+1}{n+2}I_n$であり,ウォリスの漸化式\eqref{式:ウォリス漸化式}そのものである.整数指数の多項式の積分とウォリスの積分は,置換$x=\sin^2t$によって,同じベータ積分の異なる場合として結び付いたのだ.
